\documentclass[12pt,a4paper]{article}
\usepackage[T1]{fontenc}
\usepackage[utf8]{inputenc}
\usepackage[spanish]{babel}
\usepackage{geometry}
\usepackage{booktabs}
\usepackage{xcolor}
\usepackage{listings}
\usepackage{hyperref}
\usepackage{enumitem}
\usepackage{titlesec}
\usepackage{parskip}

\geometry{margin=2.5cm}

\hypersetup{
  colorlinks=true,
  linkcolor=black,
  urlcolor=blue!70!black,
}

\lstdefinestyle{bash}{
  language=bash,
  basicstyle=\small\ttfamily,
  backgroundcolor=\color{gray!10},
  frame=single,
  rulecolor=\color{gray!40},
  breaklines=true,
  showstringspaces=false,
}

\title{%
  \textbf{Instalación de nvidia-open 580}\\[0.4em]
  \large Degradación desde 595 y compatibilidad con el kernel LTS en\\
  ASUS ROG Zephyrus G14 GA403UV con CachyOS
}
\author{pc-config}
\date{Mayo de 2026}

\begin{document}

\maketitle
\tableofcontents
\newpage

% ─────────────────────────────────────────────
\section{Contexto}

El equipo es un \textbf{ASUS ROG Zephyrus G14 GA403UV} con \textbf{CachyOS}
(base Arch Linux) y el entorno \textbf{Omarchy} (Hyprland + Wayland). Cuenta
con configuración de GPU dual:

\begin{itemize}
  \item \textbf{iGPU}: AMD Radeon RX~780M (HawkPoint) --- driver
        \texttt{amdgpu}, conectada directamente al panel.
  \item \textbf{dGPU}: NVIDIA GeForce RTX 4060 Laptop (8\,GB VRAM) ---
        driver \texttt{nvidia-open}, arquitectura PRIME Offload.
\end{itemize}

CachyOS distribuye \texttt{nvidia-open-dkms} en el repositorio
\texttt{extra}. En el momento de este procedimiento, la versión disponible
era la \textbf{595.71.05-2} (instalada inicialmente como 595.71.05-1 y
actualizada a -2). Esta versión presentaba un bug crítico de renderizado y,
además, era la única compatible con el kernel~7.x. La versión anterior,
\textbf{580.95.05-1}, no tiene el bug, pero requiere el kernel LTS~6.18.

% ─────────────────────────────────────────────
\section{Problema con nvidia-open 595: artefactos visuales}

\subsection{Síntoma}

Con \texttt{nvidia-open-dkms} 595.71.05 instalado, ciertos juegos como
\emph{The Last of Us Part~I} mostraban \textbf{artefactos visuales graves}
al ejecutarse con PRIME Offload sobre Wayland: corrupciones de textura,
franjas de píxeles incorrectos y parpadeos aleatorios que hacían los juegos
injugables.

El problema era reproducible con cualquier configuración de lanzamiento
(\texttt{prime-run}, \texttt{\_\_NV\_PRIME\_RENDER\_OFFLOAD=1}) y no
desaparecía ajustando parámetros del compositor ni del driver.

\subsection{Causa}

Se trata de un bug conocido de renderizado introducido en la serie 595 del
driver \texttt{nvidia-open} bajo Wayland con iGPU AMD en PRIME Offload. No
hay solución de configuración: la única alternativa es usar una versión del
driver sin el bug.

% ─────────────────────────────────────────────
\section{Por qué nvidia-open 580 requiere el kernel LTS}

Al intentar instalar \texttt{nvidia-open-dkms} 580 con los kernels
\textbf{7.0.8-1-cachyos} y \textbf{7.0.8-arch1-1} presentes en el sistema,
el módulo DKMS falla durante la compilación con dos errores de API del
kernel:

\begin{lstlisting}[style=bash]
# /var/lib/dkms/nvidia/580.95.05/build/make.log
error: cannot assign to non-static data member 'vm_flags' with
       const-qualified type 'const vm_flags_t'
error: call to undeclared function 'phys_to_dma'
\end{lstlisting}

\begin{itemize}
  \item \textbf{\texttt{vm\_flags}}: en kernel 7.x se cambió a un campo de
        solo lectura (\texttt{const vm\_flags\_t}). El driver 580 lo
        modifica directamente, lo que ya no está permitido.
  \item \textbf{\texttt{phys\_to\_dma}}: función eliminada o movida en
        kernel 7.x; el driver 580 la llama sin declaración previa.
\end{itemize}

DKMS registra el intento en el log de pacman:

\begin{lstlisting}[style=bash]
==> dkms install --no-depmod nvidia/580.95.05 -k 7.0.8-arch1-1
==> WARNING: `dkms install --no-depmod nvidia/580.95.05
             -k 7.0.8-arch1-1' exited 10
\end{lstlisting}

En cambio, con \textbf{linux-cachyos-lts} (6.18.x) la compilación termina
correctamente y el módulo queda instalado en:

\begin{lstlisting}[style=bash]
/lib/modules/6.18.31-1-cachyos-lts/updates/dkms/nvidia.ko.zst
\end{lstlisting}

La solución pasa por arrancar exclusivamente en el kernel LTS.

% ─────────────────────────────────────────────
\section{Obtención de los paquetes 580}

La versión 580 ya no está disponible en los repositorios activos de CachyOS.
Los paquetes deben descargarse del \textbf{Arch Linux Archive}, que conserva
versiones históricas de todos los paquetes del repositorio \texttt{extra}:

\begin{center}
\url{https://archive.archlinux.org/packages/}
\end{center}

Los paquetes necesarios y sus rutas en el archivo son:

\begin{center}
\begin{tabular}{ll}
  \toprule
  Paquete & Ruta en el archivo \\
  \midrule
  \texttt{nvidia-open-dkms} &
    \texttt{/n/nvidia-open-dkms/nvidia-open-dkms-580.95.05-1-x86\_64.pkg.tar.zst} \\
  \texttt{nvidia-utils} &
    \texttt{/n/nvidia-utils/nvidia-utils-580.95.05-1-x86\_64.pkg.tar.zst} \\
  \texttt{lib32-nvidia-utils} &
    \texttt{/l/lib32-nvidia-utils/lib32-nvidia-utils-580.95.05-1-x86\_64.pkg.tar.zst} \\
  \texttt{nvidia-settings} &
    \texttt{/n/nvidia-settings/nvidia-settings-580.95.05-1-x86\_64.pkg.tar.zst} \\
  \texttt{opencl-nvidia} &
    \texttt{/o/opencl-nvidia/opencl-nvidia-580.95.05-1-x86\_64.pkg.tar.zst} \\
  \bottomrule
\end{tabular}
\end{center}

Descargar los cinco paquetes en un directorio de trabajo:

\begin{lstlisting}[style=bash]
mkdir ~/nvidia-580 && cd ~/nvidia-580
BASE="https://archive.archlinux.org/packages"

curl -LO "$BASE/n/nvidia-open-dkms/nvidia-open-dkms-580.95.05-1-x86_64.pkg.tar.zst"
curl -LO "$BASE/n/nvidia-utils/nvidia-utils-580.95.05-1-x86_64.pkg.tar.zst"
curl -LO "$BASE/l/lib32-nvidia-utils/lib32-nvidia-utils-580.95.05-1-x86_64.pkg.tar.zst"
curl -LO "$BASE/n/nvidia-settings/nvidia-settings-580.95.05-1-x86_64.pkg.tar.zst"
curl -LO "$BASE/o/opencl-nvidia/opencl-nvidia-580.95.05-1-x86_64.pkg.tar.zst"
\end{lstlisting}

% ─────────────────────────────────────────────
\section{Procedimiento de instalación}

\subsection{Paso previo: arrancar en el kernel LTS}

El paquete \texttt{nvidia-open-dkms} ejecuta DKMS en el
\emph{post-install hook} de pacman para todos los kernels presentes en el
sistema. Si el kernel 7.x está instalado, DKMS intentará compilar para él
y fallará (salida 10), aunque el proceso continúa para el kernel LTS. Para
evitar ruido en el log y asegurarse de que el sistema activo es el LTS:

\begin{enumerate}
  \item Asegurarse de arrancar con \texttt{linux-cachyos-lts}.
  \item Opcionalmente, desinstalar \texttt{linux-cachyos} y
        \texttt{linux-cachyos-headers} si no se van a necesitar,
        para que DKMS solo construya para el kernel LTS.
\end{enumerate}

\subsection{Instalación de los paquetes locales}

Instalar los cinco paquetes en un solo comando para que pacman resuelva
correctamente las dependencias internas entre ellos:

\begin{lstlisting}[style=bash]
cd ~/nvidia-580
sudo pacman -U --noconfirm \
    nvidia-open-dkms-580.95.05-1-x86_64.pkg.tar.zst \
    nvidia-utils-580.95.05-1-x86_64.pkg.tar.zst \
    lib32-nvidia-utils-580.95.05-1-x86_64.pkg.tar.zst \
    nvidia-settings-580.95.05-1-x86_64.pkg.tar.zst \
    opencl-nvidia-580.95.05-1-x86_64.pkg.tar.zst
\end{lstlisting}

Durante el hook de pacman se verá:

\begin{lstlisting}[style=bash]
==> dkms install --no-depmod nvidia/580.95.05 -k 6.18.31-1-cachyos-lts
\end{lstlisting}

Si el kernel~7.x sigue instalado, aparecerán también los intentos fallidos
para él. Esto es normal y no impide el funcionamiento en el LTS.

\subsection{Verificar que el módulo quedó instalado}

\begin{lstlisting}[style=bash]
sudo dkms status nvidia
# nvidia/580.95.05, 6.18.31-1-cachyos-lts, x86_64: installed

ls /lib/modules/6.18.31-1-cachyos-lts/updates/dkms/ | grep nvidia
# nvidia.ko.zst  nvidia-drm.ko.zst  nvidia-modeset.ko.zst
# nvidia-peermem.ko.zst  nvidia-uvm.ko.zst
\end{lstlisting}

% ─────────────────────────────────────────────
\section{Establecer el kernel LTS como arranque predeterminado}

CachyOS usa \textbf{Limine} como gestor de arranque. Para que
\texttt{linux-cachyos-lts} sea la entrada predeterminada:

Editar \texttt{/etc/default/limine} y poner \texttt{linux-cachyos-lts} antes que
\texttt{linux-cachyos} en \texttt{BOOT\_ORDER}. Limine arranca automáticamente
la primera entrada válida. Luego regenerar la configuración con:

\begin{lstlisting}[style=bash]
sudoedit /etc/default/limine
BOOT_ORDER="linux-cachyos-lts, linux-cachyos, *fallback, Snapshots"
sudo limine-update
\end{lstlisting}

Verificar tras reiniciar:

\begin{lstlisting}[style=bash]
uname -r
# 6.18.31-1-cachyos-lts
\end{lstlisting}

% ─────────────────────────────────────────────
\section{Congelar la versión del driver}

Para que \texttt{pacman -Syu} no vuelva a actualizar \texttt{nvidia-open}
a la versión 595 automáticamente, añadir los paquetes a \texttt{IgnorePkg}
en \texttt{/etc/pacman.conf}:

\begin{lstlisting}[style=bash]
sudo nano /etc/pacman.conf
\end{lstlisting}

Localizar la línea \texttt{IgnorePkg} (o añadirla si no existe) y dejarla
así:

\begin{lstlisting}[style=bash]
IgnorePkg = nvidia-open-dkms nvidia-utils lib32-nvidia-utils \
            nvidia-settings opencl-nvidia
\end{lstlisting}

A partir de ahora, \texttt{pacman -Syu} mostrará un aviso indicando que
estos paquetes están ignorados y no los actualizará.

% ─────────────────────────────────────────────
\section{Verificación del driver}

Cargar el módulo y comprobar que la GPU es reconocida:

\begin{lstlisting}[style=bash]
# Modulo cargado
lsmod | grep "^nvidia "

# GPU detectada por el driver
nvidia-smi --query-gpu=name,driver_version --format=csv,noheader
# NVIDIA GeForce RTX 4060 Laptop GPU, 580.95.05

# Verificar PRIME Offload
prime-run glxinfo | grep "OpenGL renderer"
# OpenGL renderer string: NVIDIA GeForce RTX 4060 Laptop GPU/PCIe/SSE2
\end{lstlisting}

Para confirmar que los artefactos han desaparecido, lanzar el juego
afectado con:

\begin{lstlisting}[style=bash]
prime-run %command%
\end{lstlisting}

% ─────────────────────────────────────────────
\section{Estado final de paquetes}

Tras el procedimiento, el estado de los paquetes NVIDIA es el siguiente:

\begin{center}
\begin{tabular}{lll}
  \toprule
  Paquete & Versión instalada & Nota \\
  \midrule
  \texttt{nvidia-open-dkms}   & 580.95.05-1 & Driver principal (DKMS) \\
  \texttt{nvidia-utils}       & 580.95.05-1 & Utilidades y librerías \\
  \texttt{lib32-nvidia-utils} & 580.95.05-1 & Soporte 32\,bits \\
  \texttt{nvidia-settings}    & 580.95.05-1 & Panel de configuración \\
  \texttt{opencl-nvidia}      & 580.95.05-1 & OpenCL 64\,bits \\
  \texttt{lib32-opencl-nvidia}& 595.71.05-1 & OpenCL 32\,bits (sin degradar) \\
  \texttt{libva-nvidia-driver}& 0.0.17-1    & Aceleración de vídeo VA-API \\
  \bottomrule
\end{tabular}
\end{center}

El paquete \texttt{lib32-opencl-nvidia} quedó en la versión 595 por no haber
sido incluido en el comando de degradación. Para los casos de uso habituales
(juegos, renderizado, Wayland) esto no tiene impacto funcional.

% ─────────────────────────────────────────────
\section{Notas de mantenimiento}

\subsection{Cuándo actualizar el driver}

Antes de quitar cualquier paquete de \texttt{IgnorePkg} y actualizar a una
versión más reciente de \texttt{nvidia-open}, verificar en los foros de
CachyOS o en los \emph{release notes} de NVIDIA que:

\begin{enumerate}
  \item El bug de artefactos bajo Wayland/PRIME ha sido corregido.
  \item La nueva versión compila contra el kernel LTS activo (6.18.x) o,
        en su defecto, contra el nuevo kernel que se use.
\end{enumerate}

\subsection{Si se actualiza el kernel LTS}

DKMS recompilará automáticamente \texttt{nvidia/580.95.05} para el nuevo
kernel gracias a \texttt{AUTOINSTALL="yes"} en su \texttt{dkms.conf}. Si la
nueva versión del kernel LTS introduce cambios de API incompatibles con el
driver 580, el proceso fallará y el log estará en:

\begin{lstlisting}[style=bash]
/var/lib/dkms/nvidia/580.95.05/build/make.log
\end{lstlisting}

En ese caso habrá que valorar actualizar a una versión del driver que
compile en el nuevo kernel y no tenga el bug de artefactos.

\subsection{Restaurar la versión 595}

Si en el futuro se desea volver a la versión actual del repositorio:

\begin{lstlisting}[style=bash]
# 1. Eliminar las entradas de IgnorePkg en /etc/pacman.conf
# 2. Actualizar
sudo pacman -Syu nvidia-open-dkms nvidia-utils lib32-nvidia-utils \
                 nvidia-settings opencl-nvidia lib32-opencl-nvidia
\end{lstlisting}

\end{document}
